\section{The randomized coverage and submodular boundary}
\label{app:coverage-boundary}
This appendix strengthens the four-query upper witness while preserving
the complete source, the information actions and the full causal-skip
priority comparator. It concerns the worst-input expectation of the
\emph{final purchased-set fee}, not an expected sum of batch charges.

\begin{theorem}[Four-query coverage/submodular boundary, computer-assisted]
\label{thm:coverage-four}
In the complete equivalence-partition fibre model, the smallest indexed
pair-query menu permitting $R_{\mathcal A}<R_{\mathcal P}$ is four for
normalized nonnegative monotone submodular fees, and for nonnegative
weighted coverage fees. The same minima hold with positive nonempty fees
and strict increase under every proper inclusion. A witness already has
strictly positive individual query charges and strictly positive pairwise
shared-setup charges, without a global setup charge. It is strictly
submodular on every incomparable pair of query sets. Its exact values are
\[
 R_{\mathcal A}=\frac{125}{7},\qquad
 R_{\mathcal P}=\frac{161}{9},\qquad
 R_{\mathcal P}-R_{\mathcal A}=\frac{2}{63}.
\]
\end{theorem}
\begin{proof}
Theorem~\ref{thm:random-small} excludes a gap on complete sources of at
most three queries for every normalized nonnegative monotone fee. Each
fee class in the statement is contained in that class. We give a matching
complete four-query source and exact upper and lower certificates.

\paragraph*{Complete source.}
Use records $0,\ldots,7$ and queries $q_i=(2i,2i+1)$, $0\leq i<4$,
with integer words $w=\sum_i2^ix_i$. Under label zero, the required
SAME blocks are $\{0,1,2,4,6\},\{3,5\},\{7\}$, with records $3,7$
required DIFFERENT. Under label one they are
$\{0,2,4,6\},\{1,3,7\},\{5\}$, with no DIFFERENT requirement.
In each hypothesis retain every satisfying partition. Such a partition
is a coarsening of its three prior blocks. The five possible coarsenings
leave three legal negative partitions and five positive ones. Their
distinct projections are
\[
 W_0=\{1,7,9\},\qquad W_1=\{0,4,11,15\}.
\]
These are the same source as in
Appendix~\ref{app:randomized-adaptive-priority-gap-candidate}, not selected
subsets of its worlds. Write $\mathcal W=(0,1,4,7,9,11,15)$ for input order.

\paragraph*{Positive charges and pairwise sharing.}
For $S\subseteq\{0,1,2,3\}$, let
\begin{equation}\label{eq:pairwise-coverage-C}
 C(S)=\sum_{i\in S}c_i+
       \sum_{i<j}w_{ij}{\bf1}\{S\cap\{i,j\}\ne\varnothing\},
\end{equation}
where
\begin{align*}
 (c_0,c_1,c_2,c_3)&=(3,1,1,1),\\
 (w_{01},w_{02},w_{03})&=(2,1,2),\\
 (w_{12},w_{13},w_{23})&=(4,4,4).
\end{align*}
A shared setup is charged once when either associated query is bought;
activating it reveals no additional answer. This leaves exactly the four
original information actions for both policy classes. Every summand is
a positive weighted coverage indicator. Consequently $C$ is normalized,
nonnegative, monotone and submodular. Its least nonempty fee is $8$ and
its least proper single-query marginal is $1$. For incomparable $S,T$,
\begin{equation}\label{eq:coverage-strict-slack}
 C(S)+C(T)-C(S\cup T)-C(S\cap T)
 =\sum_{\substack{i\in S\setminus T\\j\in T\setminus S}}w_{ij}\geq1.
\end{equation}
Here $w_{ij}=w_{ji}$ when necessary. Each set difference is nonempty
and every pair weight is positive. Comparable sets necessarily give
equality. Thus the asserted strictness excludes only those identities.
In mask order the fee values are
\begin{align*}
 (C(0),\ldots,C(7))&=(0,8,11,17,10,17,17,22),\\
 (C(8),\ldots,C(15))&=(11,17,18,22,17,22,20,23).
\end{align*}

\paragraph*{All-input upper certificates.}
Write $T(q,u,v)$ for buying $q$ and following $u$ on zero and $v$ on
one; constants denote labels. Set $u=T(0,1,0)$, $v=T(1,1,0)$,
$z=T(3,u,0)$. Use
\begin{align*}
 a_1&=T(1,u,T(3,0,1)),\\
 a_2&=T(3,u,T(1,0,1)),\\
 a_3&=T(2,T(1,z,1),T(3,v,1)),\\
 p_3&=T(2,T(1,u,1),T(3,v,1)),\\
 p_4&=T(2,T(1,z,1),T(3,u,1)).
\end{align*}
Their paid-set masks on $\mathcal W$ are respectively
\begin{align*}
 a_1 &: (3,3,3,10,3,10,10),\\
 a_2 &: (9,9,9,9,10,10,10),\\
 a_3 &: (15,15,14,14,14,6,12),\\
 p_3 &: (7,7,14,14,7,6,12),\\
 p_4 &: (15,15,13,13,14,6,12).
\end{align*}
All five trees are correct on all seven inputs. The four trees
$a_1,a_2,p_3,p_4$ admit global orders $(1,0,2,3)$, $(2,3,0,1)$,
$(2,3,1,0)$ and $(2,1,3,0)$, respectively, with causal skips.
Mixing $a_1,a_2,a_3$ with weights $(3/7,3/7,1/7)$ gives input costs
\[
 (125/7,125/7,122/7,125/7,125/7,125/7,125/7).
\]
Mixing $a_1,a_2,p_3,p_4$ with $(4/9,4/9,1/18,1/18)$ gives
\[
 (317/18,317/18,157/9,161/9,161/9,161/9,161/9).
\]
These prove $R_{\mathcal A}\leq125/7$ and
$R_{\mathcal P}\leq161/9$. The priority mixture is an achievable
witness, not a restriction of the comparison class to four templates.

\paragraph*{Lower certificates for the full classes.}
Use priors on $(7,9,11,15)$
\[
 \pi_A=(1/7,1/7,5/14,5/14),\qquad
 \pi_P=(1/9,1/9,7/18,7/18).
\]
The other three inputs remain correctness obligations with prior zero.
Every correct deterministic adaptive tree has $\pi_A$ mean at least
$125/7$, and every correct global-order priority tree has $\pi_P$ mean
at least $161/9$. The complete retained families have $40$ adaptive and $32$ priority
paid-set profiles; pointwise dominated profiles may remain. Their completeness
and pointwise paid-set domination are inherited from the source/family
proofs used in Section~\ref{sec:four-additive}; no new catalogue is assumed.
Direct feasible-state reconstruction gives $35$ adaptive and $27$
priority pointwise-minimal profiles, with five redundant profiles in
each retained family. Every minimal profile occurs in its retained
family, and every retained profile has a reconstructed profile whose
paid set is a subset on every input. The original legal paths and
global orders were also checked directly. All retained columns are
kept in the certificates below; their bounds are unchanged.
The reduced costs for all these columns, multiplied by $14$ and $18$
respectively, are the following concatenated rows:
\[
\begin{array}{rrrrrrrrrr}
48&55&58&48&55&65&0&9&58&44\\
33&29&65&51&40&36&33&10&19&8\\
4&0&9&2&11&11&44&30&10&19\\
6&15&29&6&15&4&2&11&11&0
\end{array}
\]
and
\[
\begin{array}{rrrrrrrr}
64&73&74&64&73&83&0&11\\
74&56&39&35&83&44&39&10\\
21&4&0&0&11&2&13&13\\
10&21&6&17&6&17&0&2.
\end{array}
\]
All entries are nonnegative. Under a monotone fee, a pointwise
paid-set-dominated profile cannot improve the lower bound. The exact
checker re-traced all $504$ original policy/input paths and their order
witnesses, and checked every displayed column inequality. Mixing private
seeds preserves a fixed-prior lower bound, and a maximum over inputs
dominates that mean. The two priors differ; this is not an assertion
of adaptive/priority equality for every specified prior.

There is also a separate full-class check that uses no retained family
or first checker result. It reconstructs the physical source above.
For a prior $\pi$, feasible input set $U$ and already purchased set $B$,
the pure-state value is $\pi(U)C(B)$. At an impure state, the adaptive
recurrence minimizes the sum of its two child contributions over every
unbought query. At a fixed priority order position it considers both
permanent SKIP and TAKE, and advances the same order in both answer
children. An impure exhausted state is illegal even if $\pi(U)=0$.
Induction on remaining queries or positions proves that these recurrences
cover every legal deterministic policy. Minimizing over all $24$ orders
covers full $\mathcal P$. Randomizing cannot reduce a fixed-prior mean
below the best deterministic value.

The two actual checks used the single frozen fee
$F(S)=C(S)+9{\bf1}\{S\ne\varnothing\}$. They obtained exact lower
and upper values $188/7$ and $242/9$. Because both labels occur, every
correct deterministic seed buys a nonempty set on every input. Removing
the global charge therefore subtracts exactly $9$ from every policy/input
cost, and from each mixed worst-input objective. This proves the claimed
$C$ values without a second cost-game optimization. The separate check
uses $63$ adaptive states and $1602$ states over the $24$ priority orders.
The arbitrary-policy lower justification is the stated recurrence
induction, supported by these exact finite calculations.
For private seeds defined only almost surely, discard the finite union
of erroneous null seed sets over the seven inputs first.
\end{proof}

\subsection{How sharing repairs the earlier fee}
Let $F_0$ be the old nonsubmodular fee in
Appendix~\ref{app:randomized-adaptive-priority-gap-candidate} and put
$H_1(S)=F_0(S)+|S|$. Directly on all $16$ sets,
\begin{align*}
 F(S)=H_1(S)&+5{\bf1}\{S\cap\{0,1,2\}\ne\varnothing\}\\
 &+6{\bf1}\{S\cap\{0,1,3\}\ne\varnothing\}\\
 &+5{\bf1}\{S\cap\{0,2,3\}\ne\varnothing\}\\
 &+8{\bf1}\{S\cap\{1,2,3\}\ne\varnothing\}.
\end{align*}
If $|S|\geq2$, all four triples are met, so $F(S)=H_1(S)+24$ and
$C(S)=H_1(S)+15$. Every input in the dual support $(7,9,11,15)$
has an opposite-label input agreeing with each possible first query
answer. Every correct policy therefore buys at least two queries there.
The five explicit upper trees buy at least two on all seven inputs.
The old matched $H_1$ bounds $20/7$ and $26/9$ consequently transfer
by adding $15$, giving another explanation of the same gap.

It is false that every original input requires two queries: inputs
$0,4$ have one-query certificates at $q_0$. They are not deleted. The
argument needs two-query necessity only on the dual support, plus the
all-input property of the displayed upper trees. In particular it makes
no uniform $H_1$ cost-shift claim for arbitrary policies on the other
inputs. The uniform deletion of the \emph{global} charge $9$, in contrast,
uses only nonempty purchase sets and is valid for all correct policies.

\begin{corollary}[Robust submodular and coverage neighborhoods]
\label{cor:coverage-robust}
For the fixed source and fee $C$ in Theorem~\ref{thm:coverage-four}, let
$G(\varnothing)=0$ and $\max_{S\ne\varnothing}|G(S)-C(S)|\leq1/126$.
Every such $G$ is positive on nonempty sets, strictly increasing, and
strictly submodular on incomparable pairs, with
\[
 R_{\mathcal P}(G)-R_{\mathcal A}(G)\geq\frac1{63}.
\]
There is also a coverage neighborhood: perturb each of the ten positive
parameters in~\eqref{eq:pairwise-coverage-C} by at most $1/1260$.
Every resulting fee has these properties and the same gap lower bound.
\end{corollary}
\begin{proof}
If the per-set fee change is at most $r$, each policy's cost on each
input changes by at most $r$. The explicit adaptive mixture and full
priority dual from Theorem~\ref{thm:coverage-four} give
\[
 R_{\mathcal P}(G)-R_{\mathcal A}(G)\geq2/63-2r.
\]
At $r=1/126$ this is $1/63$. Nonempty fees remain at least $8-r>0$;
proper single-query marginals remain at least $1-2r>0$; incomparable
submodular slacks remain at least $1-4r>0$. Comparable inequalities
are identities. Thus the full $15$-dimensional normalized real-fee box
lies inside the specified submodular class. It need not consist of
coverage functions. For the ten-parameter box, every parameter remains
positive and every set fee changes by at most $10/1260=1/126$.
The previous argument applies. These radii are conservative analytic
guarantees, not maximum-radius or sampled-neighborhood claims.
\end{proof}

\subsection{Evidence, prior costs and remaining scope}
The main rational check uses the original $40/32$ families and checks
all $256$ submodular inequalities and $64$ query-addition conditions,
including the $32$ proper additions. A separate implementation uses the
physical constraints and frozen priors rather than those policy tables.
Both actual processes exited zero. They were written and run by the
same AI agent; algorithmic separation is not human or external review.
The initial concave-cardinality repair candidate was retained unexecuted.
The coverage fee and certificates were fixed before their first check;
global-charge removal and the neighborhood bounds are analytic corollaries.
No new fee/prior search, LP solver, model API or hosted run was used.

Coverage costs and shared resources are established models. In particular,
Harris et al.~\cite{harris2025} include coverage machine costs in their
batch-testing setting. Their independent series-diagnosis objective
charges the sum of performed batch costs in expectation. Here the
source is a correlated complete label promise, the objective maximizes
over inputs after private expectation, and each path pays its final
union fee once. These differences require a cost-, source-, objective-
and comparator-preserving reduction before transferring a theorem.
No such subsuming reduction is established here, and arbitrary alternative
reductions are not ruled out. Exact bibliographic priority remains open.

The minimum four does not classify all four-query submodular tasks.
The $1345$ source orbits already equal for all monotone fees remain equal
for submodular fees. The witness orbit has $48$ physical coordinate/label
transports; this follows by symmetry, not $48$ new runs. Ten other
four-query source orbits remain unclassified for general submodular fees.
The additive exclusion through four and the interval
$5\leq m_{\rm add}\leq m_{\rm pos}\leq7$ are unchanged. The displayed
ratio $1127/1125$ is one exact witness value, not a largest four-query
ratio. No natural academic-evaluation advantage, new human gold labels,
human proof review, formal-kernel certification or historical execution
authentication follows from this theorem.
